Substance use disorders are an important determinant of maternal and infant health outcomes in the United States \citep{jarlenski2020substance, clemans2019pregnant}. Identifying maternal drug use during pregnancy may improve maternal and infant health outcomes through pharmacological treatments, counseling, and postpartum services. However, infant and maternal drug tests are not used exclusively for medical decision-making. Drug test results also are sometimes reported to child protective services and criminal justice agencies. Although identifying substance use and child neglect are relevant to the missions of child protective services and criminal justice agencies, research suggests some pregnant women fear that a positive drug test puts them at risk for child separation and criminal justice system involvement \citep{roberts2011complex}. These fears may lead people to delay or avoid prenatal care, attempt a home birth, or attempt to detox without the supervision of a doctor \citep{stone2015pregnant, slaughter2019skin}.

There is limited evidence on how often infant and maternal drug tests lead to non-medical consequences and on how these patterns differ by racial group. \citet{mccabe2022criminalization} conducts qualitative interviews with health care providers and argues that providers seem to order tests for non-therapeutic reasons related to criminal justice concerns. In a retrospective chart review at a health system in Massachusetts, \citet{cohen2023disparities} found that about 4.8\% of Black infants and 1.8\% of White infants were tested for in-utero drug exposure. Among the infants who were drug tested, there was no racial disparity in the fraction who were referred to social work assessment or child protective services, or who experienced non-parental custody placement at discharge. Beyond the domain of infant and maternal health, research on physician expertise, predictive algorithms, and multistage assessment systems also suggests that there are racial disparities in medical treatment patterns that cannot be justified on purely clinical grounds \citep{currie2017diagnosing, obermeyer2019dissecting, baron2024discrimination}.

Maternal and infant drug-testing practices are not standardized, and procedures for reporting drug tests to child protective agencies vary across states and across hospitals within states \citep{national2018understanding, mascola2017opioid}. The American College of Obstetricians and Gynecologists (ACOG) recommends infant drug tests at birth only when physicians suspect exposure to harmful substances and a test is necessary for treatment. ACOG guidelines encourage physicians to assess an infant's risk of in-utero drug exposure in terms of the infant's health state (e.g. low birthweight or symptoms of withdrawal), and whether the mother's clinical record includes a history of drug misuse, late or inadequate prenatal care, and/or current tobacco or alcohol use \citep{palmer2017evaluating}. Most previous work on racial disparities in infant and maternal drug testing is based on relatively small samples collected from a single medical center or hospital network \citep{simpson2022prediction,Soos2024,schoneich2023incidence,kunins2007effect}. This work suggests Black infants are tested at higher rates than White infants. However, existing research does not measure the degree to which racial differences in infant drug testing rates reflect differences in the clinical risk factors that are supposed to guide physician testing decisions.  

In this study, we examine electronic medical records (EMR) on a large sample of $118,193$ White and $22,369$ Black newborn infants who were born between 2019 and 2023 at facilities that belong to the Indiana Network for Patient Care, which includes most of the hospitals and clinics in Indiana. For each infant, we observe whether a meconium or umbilical cord drug test was conducted and whether the test indicated in-utero drug exposure.\footnote{Meconium and umbilical cord infant drug tests measure maternal drug use in the second and third trimester of pregnancy \citep{nms_labs_newborn_2021}.} We link the records of infants and mothers and use them to measure clinically relevant risk factors.  

We find that 16\% of Black infants and 7\% of White infants were tested. Among the infants who were tested, 20\% of Black infants tested positive and 20\% of White infants tested positive. At first glance, equal positive testing rates may suggest racial disparities in testing reflect differences in underlying risk rather than a bias towards testing Black infants. However, identical positive testing rates are also consistent with a lower testing threshold for Black infants; physicians might over test low-risk Black infants, bringing down average positive testing rates. Without knowledge of the distribution of risk factors in the Black and White tested populations, these average positive testing rates alone do not reveal whether testing decisions involve racial prejudice or animus \citep{becker1957economics, knowles2001racial, ayres2002outcome}. The EMR data show that White and Black infants do have a different distribution of clinical risk factors. For example, mothers of Black infants had higher prevalence of maternal substance use history and resided in counties with higher rates of neonatal abstinence syndrome and drug overdose mortality. We use inverse propensity score weights to decompose the racial testing gap into a component explained by differences in clinically relevant risk factors and an unexplained component, which may reflect racial bias. The measured risk factors in our data account for approximately $39\%$ of the Black-White testing gap. This implies that about $61\%$ of the gap in testing arises because of differential treatment of infants with the same risk, or because clinicians may observe additional risk factors that are not captured in our data. 

